\setcounter{section}{3}
\section{Memoria de cálculo eléctrico explícita}
\setcounter{subsection}{0}
\subsection{Datos, resistencia y reactancia}
Conductor propuesto AAAC 6201-T81,50mm\textsuperscript{2} nominales,49.5mm\textsuperscript{2} efectivos, diámetro 9 mm, masa 135 kg/km, R20=0.671 ohm/km. Masa y diámetro concuerdan con la tabla de referencia DGE; resistencia, RTS 14.7 kN, E62 GPa y curva de fluencia siguen como datos tipo sin fabricante contratado.

R75 = R20\ensuremath{\times}[1+\ensuremath{\alpha}R\ensuremath{\times}(75-20)]\ensuremath{\times}Lc/1000 =0.671\ensuremath{\times}(1+0.0036\ensuremath{\times}55)\ensuremath{\times}6.758120 = 5.432569 ohm/fase. Se usa longitud física de conductor; R01 usaba proyección horizontal y subestimaba ligeramente R.

DMG=(Dab\ensuremath{\times}Dbc\ensuremath{\times}Dca)\textasciicircum{}(1/3)=(1.2\ensuremath{\times}1.2\ensuremath{\times}2.4)\textasciicircum{}(1/3)=1.511905 m. GMR\ensuremath{\approx}0.7788\ensuremath{\times}d/2=0.0035046 m, aproximación a sustituir por GMR de cable trenzado. X=2\ensuremath{\pi}f\ensuremath{\times}2\ensuremath{\times}10\textasciicircum{}-4\ensuremath{\times}ln(DMG/GMR)\ensuremath{\times}Lc(km)=3.091467 ohm/fase, a 60 Hz.

\setcounter{subsection}{1}
\subsection{Qué significa que la salida sea 10 kV}
Primera frontera: si se fijan exactamente 10 kV entre fases en el extremo MT de salida, se calcula sólo línea y carga. Segunda frontera: si se fijan 380 V nominales en la barra equivalente del generador, T1 y T2 agregan sus propias impedancias y pérdidas; la MT inicial ya no es exactamente 10 kV. Ambas fronteras se tabulan por separado para evitar respuestas incompatibles.

\begingroup\small\setlength{\LTpre}{4pt}\setlength{\LTpost}{6pt}
\begin{longtable}{*{4}{L{\dimexpr\linewidth/4-2\tabcolsep\relax}}}
\toprule
\textbf{MT salida fija kV} & \textbf{MT llegada kV} & \textbf{Corriente A} & \textbf{Pérdida kW} \\ \midrule
\endfirsthead
\toprule
\textbf{MT salida fija kV} & \textbf{MT llegada kV} & \textbf{Corriente A} & \textbf{Pérdida kW} \\ \midrule
\endhead
10 & 9.8797 & 11.135 & 2.0206 \\
20 & 19.9404 & 5.517 & 0.4960 \\
22 & 21.9458 & 5.013 & 0.4095 \\
22.9 & 22.8480 & 4.815 & 0.3778 \\
\bottomrule\end{longtable}\endgroup

\setcounter{subsection}{2}
\subsection{Flujo con ambos transformadores}
Base 250 kVA. rTX=Pcu/Sn=3.1/250=0.0124 pu; xTX=sqrt(0.06\textsuperscript{2}-0.0124\textsuperscript{2})=0.058705 pu. Zbase=Vnom\textsuperscript{2}/Sbase; ZTX=(r+jx)\ensuremath{\times}Zbase. P0=0.45 kW y Q0=2.5kvar por TX son supuestos constantes. Los equipos alternativos se comparan con las mismas pérdidas pu; ofertas reales pueden cambiar el resultado.

Barrido por fase: i\_\allowbreak{}carga=conj[S\_\allowbreak{}trifásica/(3v)]; corriente de rama=suma de cargas aguas abajo; v\_\allowbreak{}receptor=v\_\allowbreak{}emisor-Z\ensuremath{\times}i; iterar hasta residuo<10\textasciicircum{}-8 V. P\_\allowbreak{}Joule=3|i|\textsuperscript{2}R; Q\_\allowbreak{}serie=3|i|\textsuperscript{2}X. Comprobar P\_\allowbreak{}fuente=P\_\allowbreak{}carga+P0\_\allowbreak{}T1+P0\_\allowbreak{}T2+\ensuremath{\Sigma}P\_\allowbreak{}Joule y Q equivalente. Añadir5 kW de auxiliares sólo al balance de generación; su FP se supone 1.

\begingroup\small\setlength{\LTpre}{4pt}\setlength{\LTpost}{6pt}
\begin{longtable}{*{6}{L{\dimexpr\linewidth/6-2\tabcolsep\relax}}}
\toprule
\textbf{kV nominal} & \textbf{Fuente pu} & \textbf{MT salida/llegada kV} & \textbf{BT V} & \textbf{P línea kW} & \textbf{P generación kW} \\ \midrule
\endfirsthead
\toprule
\textbf{kV nominal} & \textbf{Fuente pu} & \textbf{MT salida/llegada kV} & \textbf{BT V} & \textbf{P línea kW} & \textbf{P generación kW} \\ \midrule
\endhead
10 & 1.0 & 9.660 / 9.530 & 350.3 & 2.359 & 184.05 \\
10 & 1.025 & 9.921 / 9.795 & 360.7 & 2.224 & 183.67 \\
10 & 1.05 & 10.181 / 10.059 & 371.1 & 2.102 & 183.33 \\
20 & 1.0 & 19.329 / 19.265 & 354.3 & 0.576 & 182.17 \\
20 & 1.025 & 19.850 / 19.788 & 364.6 & 0.544 & 181.91 \\
20 & 1.05 & 20.370 / 20.309 & 374.8 & 0.515 & 181.66 \\
22 & 1.0 & 21.263 / 21.204 & 354.5 & 0.476 & 182.07 \\
22 & 1.025 & 21.836 / 21.779 & 364.8 & 0.449 & 181.81 \\
22 & 1.05 & 22.407 / 22.352 & 375.0 & 0.425 & 181.57 \\
22.9 & 1.0 & 22.133 / 22.076 & 354.6 & 0.439 & 182.03 \\
22.9 & 1.025 & 22.729 / 22.675 & 364.9 & 0.414 & 181.77 \\
22.9 & 1.05 & 23.324 / 23.271 & 375.1 & 0.392 & 181.53 \\
\bottomrule\end{longtable}\endgroup

La elevación efectiva 1.025/1.05 pu es una sensibilidad del modelo ideal; no ordena cambiar un tap. Comprobar relación real, sentido del selector, impedancias por toma, máximo en vacío, rango AVR y garantía a la altitud. Debe evaluarse carga mínima y máxima y toda la red posterior. El objetivo\ensuremath{\pm}5\% de estudio no constituye certificación de calidad en medidores.

\setcounter{subsection}{3}
\subsection{Comparación económica y necesidad de regulación}
A 1.05 pu,10 kV pierde 2.102 kW en línea y 22.9 kV pierde 0.392 kW: ahorro de pico 1.710 kW. La mejora de BT también depende de T1/T2. No aparece justificación de un elevador intermedio sólo por 6.7km: primero resolver tensión nominal, transformadores, AVR y perfil de carga. La inversión adicional no puede decidirse por menor pérdida sin precio de equipos y energía.

Energía anual de pérdidas =\ensuremath{\Sigma}(P\_\allowbreak{}perdida\_\allowbreak{}i\ensuremath{\times}horas\_\allowbreak{}i), con 8760horas repartidas 25/50/75/100\% del pico. Esta curva es sintética. La hoja energia.csv conserva cada bloque. Recuperación simple de sobrecosto: \ensuremath{\Delta}CAPEX/(\ensuremath{\Delta}E\_\allowbreak{}kWh\ensuremath{\times}costo\_\allowbreak{}energía-\ensuremath{\Delta}OPEX); faltan cotizaciones y curva horaria para dar un plazo válido.

\setcounter{subsection}{4}
\subsection{Cortocircuito, protección y puesta a tierra}
Ik3\ensuremath{\approx}VLL/[sqrt(3)|Zg+ZT1+Zlínea|]; se distinguen Xd\ensuremath{^{\prime\prime}}0.20 pu, Xd\ensuremath{^{\prime}}0.30 pu y Xd1.80 pu de estudio. La corriente de falla sostenida puede ser menor que la corriente de carga: no basta elegir fusible o50/51 copiando un proyecto SEIN. Requiere curva de decremento/AVR, arranque de motores, inrush, TC/TT, lógica 51V u otra protección y coordinación real. El CSV no es ajuste de relés.

Si se mantiene MT delta en ambos TX, no hay neutro MT accesible; NGR exige transformador de puesta a tierra. Para 5 A: RN=VLL/(sqrt3\ensuremath{\times}5); a 10 kV=1154.7 ohm y 22.9 kV=2644.3 ohm. Energía nominal 10s=Vfase\ensuremath{\times}5\ensuremath{\times}10; dimensionar equipos y verificar régimen, capacitancias y protección. No sumar un NGR simbólico sin definir el neutro.

La resistencia de tierra, tensiones de paso/contacto y cimentaciones quedan pendientes. Se requieren Wenner, corriente/tiempo de falla y modelo del terreno. La ampacidad del conductor requiere datos de proveedor y balance térmico a la altitud. Una corrida con h supuesto o el bajo valor de corriente no acredita por sí sola capacidad garantizada.

